%\documentclass{amsart}
\documentclass[10pt,a4paper,leqno]{article}
%%%%%%%%%%%%%%%%%%%%%%%packages%%%%%%%%%%%%%%%%%%%%%
\usepackage[T1]{fontenc}
\usepackage[latin1]{inputenc}
\usepackage{float}
%\usepackage[notref,notcite]{showkeys}
\usepackage{graphics}
\usepackage{amssymb}
\usepackage{float}
\usepackage{makeidx}
\usepackage{graphicx}
\usepackage{amsmath} 
\usepackage{pstricks}

%%%%%%%%%%%%%%%%%%mise en page %%%%%%%%%%%%%

\textwidth17cm \textheight23truecm
\hoffset0truecm \voffset-1truecm 
\parindent 0truemm 
\evensidemargin0cm
\oddsidemargin 0cm
\topmargin1cm
\parskip-0.05cm
\renewcommand{\baselinestretch}{1.05} 
\catcode`\�=\active \def �{\' e}

 
%%%%%%%%%%%%%%%%%%%%%%%fontes (� v�rifier) %%%%%%%%%%%%%%%%%
\newfont{\Blackboard}{msbm10 scaled 1200}
\newcommand{\bll}[1]{\mbox{\Blackboard #1}}
\newfont{\roma}{cmr10 scaled 1200}
 
%%%%%%%%%%%%%%%%%%%%%%%%%%%%%%% Les ensembles de nombres %%%%

\def \Z {{\mathbb{Z}}}
\def \R {{\mathbb{R}}}
\def \N {{\mathbb{N}}}
\def \C {{\mathbb{C}}}
\def \Q {{\mathbb{Q}}}
%%%%%%%%%%%%%%%%%%commandes usuelles 

\newcommand{\dint}{\displaystyle\int}
\newcommand {\nc}   {\newcommand}
\newcommand{\dsum}{\displaystyle \sum}
\newcommand{\dps}{\displaystyle }
\def\dfrac{\displaystyle \frac }
\def\edc{\end{document}}

%%%%%%%%%%%%commandes usuelles bis 

\nc {\be}   {\begin{equation}} \nc {\ee}   {\end{equation}} 
\nc {\beq}  {\begin{eqnarray}} \nc {\eeq}  {\end{eqnarray}}
\nc {\beqs} {\begin{eqnarray*}} \nc {\eeqs} {\end{eqnarray*}}
\newcommand{\bthe}{\begin{theorem}}
\newcommand{\ethe}{\end{theorem}}
\newcommand{\brk}{\begin{remark}}
\newcommand{\erk}{\end{remark}}
\newcommand{\bco}{\begin{corollary}}
\newcommand{\eco}{\end{corollary}}
\newcommand{\blem}{\begin{lemma}}
\newcommand{\elem}{\end{lemma}}
\newcommand{\bprop}{\begin{proposition}}
\newcommand{\eprop}{\end{proposition}}
\newcommand{\bex}{\begin{example}}
\newcommand{\eex}{\end{example}}
%\newcommand{\proof}{{\bf Proof: }}
\newcommand{\fin}{\,\rule{1ex}{1ex}\,}
\newcommand {\proof} {\noindent \bf Proof. \rm}
%\newcommand{\findem}{{\unskip \nobreak \hfil \penalty 50\hskip 2em\hbox
%{}\nobreak



%%%%%%d�finition modifiable selon l'article%%%%%%
\def \eps {{\epsilon}}
\def \ga {{\gamma}}
\def \om {{\omega}}
\def \lb {{\lambda}}
\def \het {{\hat{\eta}}}
\def \hf {{\hat{f}}}
\def \hy{{\hat{y}}}

\newcommand{\caS}{{\cal S}}
\newcommand{\caD}{{\cal D}}
\newcommand{\caN}{{\cal N}}
\newcommand{\caM}{{\cal M}}
\newcommand{\caA}{{\cal A}}
\newcommand{\caI}{{\cal I}}
\newcommand{\caL}{{\cal L}}
\newcommand{\caR}{{\cal R}}
\newcommand{\caW}{{\cal W}}
\newcommand{\caV}{{\cal V}}
\newcommand{\caH}{{\cal H}}
\newcommand{\caP}{{\cal P}}
\newcommand{\caF}{{\cal F}}
\newcommand{\caC}{{\cal C}}
%\newcommand{\lb}{\lambda}
\newcommand{\da}{\caD(\caA)}
\newcommand{\ds}{\displaystyle}



\newcommand{\norm}[2]{\|#1 \| _{#2} }

% Version du 20 juillet 2023



%%%%%%%%%%%%%%%%%%%%def


\newtheorem{theorem}{Theorem}[section]
\newtheorem{lemma}[theorem]{Lemma}
\newtheorem{corollary}[theorem]{Corollary}
\newtheorem{remark}[theorem]{Remark}
\newtheorem{definition}[theorem]{Definition}
\newtheorem{proposition}[theorem]{Proposition}
\newtheorem{example}[theorem]{Example}
%%%%%%%%%%%%%%%%%%%%%%%%%%%%%%%
%\title{Asymptotic spectrum of a Bresse system}

%%  38 caract�res (avec les espaces) 

\begin{document}
\thispagestyle{empty}
\title{\bf The lack of exponential stability of a Bresse system subjected only to two dampings}
\author{Virginie R\'egnier \thanks{Univ. Polytechnique Hauts de France, INSA Hauts de France, CERAMATHS-Laboratoire de Mat�riaux C�ramiques et de Math�matiques, F-59313 Valenciennes, France, e-mail: virginie.regnier@uphf.fr} \; Wael Youssef \thanks{Department of Mathematics, Faculty of Sciences 1, Lebanese University, Hadath, Beyrouth, Lebanon, e-mail: Youssef.Wael@ul.edu.lb}}
\date{}
\maketitle


%%%%%%%%%%%%%%%%%%%%%%%%%%%%%%%%%%%%%%%%%%%%%%%%%%%%%%%%%%%%%%%%%%%


%% Plan du papier


%%%%%%%%%%%%%%%%%%%%%%%%%%%%%%%%%%%%%%%%%%%%%%%%%%%%%%%%%%%%%%%%%%%


\begin{abstract}
In this paper, we study the indirect boundary stabilization of a Bresse system with only two dissipation laws. This system, which models the dynamics of a beam, is a hyperbolic system with three wave speeds. We study the asymptotic behaviour of the eigenvalues and of the eigenvectors of the underlying operator in the case of three distinct wave velocities which is not physically relevant. Since the imaginary axis is proved to be an asymptote for one family of eigenvalues, the stability can not be exponential. Of course, this paper is only interesting from a mathematical point of view.  
\end{abstract}


\noindent \bf Key words: \rm Bresse beam; different wave propagation speeds; eigenvalues; asymptotic behaviour; non-exponential stability. \\
%indirect boundary damping
\\ 
\bf Mathematics Subject Classification: \rm 74K10, 35B40, 35Q74, 35B35, 35E15, 34L10, 34L15, 34L20, 93D15  

\maketitle
%%%%%%%%%%%%%%%%%%%%%%%%%%%%%%%%%%%%%%%%%%%%%%%%%%%%%%%%%%%%%%%%%%%%%%%%%%%%%%%%%%%%%%%%%%%%
%%%%%%%%%%%%%%%%%%%%%%%%%%%%%%%%%%%%%%%%%%%%%%%%%%%%%%%%%%%%%%%%%%%%%%%%%%%%%%%%%%%%%%%%%%%%

%\section{Conflict of interest} 

%Not applicable.

\section{Introduction}\label{intro}
\setcounter{equation}{0}

\noindent We consider the following Bresse-type system:\\

\begin{equation}
\left\{\begin{array}{ll}
\rho_1 \tilde{\varphi}_{tt}- \kappa (\tilde{\varphi}_x+ \tilde{\psi} +l \tilde{\omega})_x-\kappa_0 l(\tilde{\omega}_x-l \tilde{\varphi})=0\quad\text{in }(0,1)\times (0,\infty)\\ 
\rho_2 \tilde{\psi}_{tt}- \tilde{b} \tilde{\psi}_{xx}+\kappa(\tilde{\varphi}_x+ \tilde{\psi} +l \tilde{\omega})= 0\quad\text{in }(0,1)\times (0,\infty)\\ 
\rho_1 \tilde{\omega}_{tt}-\kappa_0 (\tilde{\omega}_x-l \tilde{\varphi})_x+\kappa l (\tilde{\varphi}_x+ \tilde{\psi} + l \tilde{\omega})=0\quad\text{in }(0,1)\times (0,\infty) 
\label{e11}
\end{array}\right.
\end{equation}
\\
with $t>0$ and $0<x<1$.\\ 
It is also called circular arch problem. The length of the beam here is $1$. The functions $\tilde{\varphi}$, $\tilde{\psi}$ and $\tilde{\omega}$ represent the longitudinal, vertical and shear angle displacements of the beam respectively. \\
The physical constants are: $\rho_1 = \rho A$, $\rho_2 = \rho I$, $\kappa = \kappa' GA$, $\kappa_0 = EA$, $\tilde{b} = EI$, $l = R^{-1}$, where $\rho$ is the density of the material, $E$ is Young's modulus, $G$ is the shear modulus, $\kappa'$  is the
shear factor, $A$ is the cross-sectional area, $I$ is the second moment of area of the cross-section and $R$ is the radius of curvature. \\ 
The Bresse system is a generalization of the Timoshenko system which is obtained by letting $R$ tend to $\infty$, or $l$ to $0$. \\
To get a dimensionless problem, as in our paper with D. Mercier (\cite{EECT18}), we consider the change in the variable $t$: $\varphi(x,t)= \tilde{\varphi} \left( x,\sqrt{\dfrac{\rho_1}{\kappa}} t \right)$, $\psi(x,t)= \tilde{\psi} \left( x,\sqrt{\dfrac{\rho_1}{\kappa}} t \right)$, $\omega(x,t)= \tilde{\omega} \left( x,\sqrt{\dfrac{\rho_1}{\kappa}} t \right)$, $a =\dfrac{\kappa_0}{\kappa }$, $b= \dfrac{\tilde{b} \rho_1}{\kappa \rho_2}$, $c = \dfrac{\rho_1}{\rho_2}$. The same effort to explain the dimensionless expression of the problem is done in the paper by L. Zietsman, N.F.J. van Rensburg and A.J. van der Merwe (\cite{Z-R-M}) for a Timoshenko system with slightly different notation. \\
\\
\noindent Finally, the system we consider in the following is:

\begin{equation}
\left\{\begin{array}{llll}
\varphi_{tt} - (\varphi_x + \psi + l \omega)_x - a l (\omega_x - l \varphi) &=& 0  \hskip 0.5cm \mbox{in } (0,1) \times (0,\infty) \\
\psi_{tt} - b \psi_{xx} + c(\varphi_x + \psi + l \omega)&=&0 \hskip 0.5cm \mbox{in } (0,1) \times (0,\infty) \\
\omega_{tt}- a (\omega_x - l \varphi)_x + l(\varphi_x + \psi + l \omega)&=&0 \hskip 0.5cm \mbox{in } (0,1) \times (0,\infty)
%u_x(0,t)=y(0,t)=u(1,t)&=&0  \hskip 0.5cm \mbox{in } (0,\infty)
\label{bresse}
\end{array} \right.
\end{equation}

\noindent with the initial conditions: for $x \in (0;1)$

\begin{equation} 
\left\{ \begin{split}
& \varphi(x,0)= \varphi_0(x), \ \varphi_t(x,0)=\phi_1(x), \ \psi(x,0)=\psi_0(x), \ \psi_t(x,0)=\psi_1(x) \\
&  \omega(x,0)= \omega_0(x), \ \omega_t(x,0)= \omega_1(x)
\label{IC}
\end{split} \right.
\end{equation}

\noindent where $a$, $b$, $c$ and $l$ are strictly positive constants. \\
\\
\noindent We consider the boundary conditions:

\begin{equation}
\left \{ \begin{array}{ll}
\varphi(t,0)=\varphi(t,1)=\psi(t,0)=\omega(t,0)=0 \\
b c^{-1} \psi_x(t,1)= -\beta_1 \psi_t(t,1),\, a(\omega_x-l\varphi)(t,1)= -\beta_2 \omega_t(t,1)
 \label{BC}
\end{array} \right.
\end{equation}

\noindent where $\beta_1$ and $\beta_2$ are strictly positive constants.\\
\\
%%%%%%%%%%%%%%%%%%%%%%%%%%%%%%%%%%%%%%%%%%%%%%%%%%%%%%%%%%%%%%%%

\noindent The energy of the solutions of System (\ref{bresse}) is defined by
\begin{equation}
E(t)=\dfrac{1}{2} \displaystyle\int_{0}^1 \Big(\vert \varphi_t \vert^2 + \dfrac{1}{c} \vert \psi_t \vert^2 + \vert \omega_t \vert^2 + \dfrac{b}{c} \vert \psi_x \vert^2 + \vert \varphi_x + \psi + l \omega \vert^2 + a \vert \omega_x-l \varphi \vert^2 \Big)dx.
\label{def.energy}
\end{equation}

\noindent Moreover, the dissipation relation is

\begin{equation}
\dfrac{dE}{dt}(t)= -\beta_1 \vert\psi_t(t,1)\vert^2 - \beta_2 \vert \omega_t(t,1) \vert^2 \leq 0.
\label{dissipation}
\end{equation} 

\noindent Hence System (\ref{bresse}) is dissipative in the sense that its associated energy is non increasing with respect to time. \\
\\
\noindent Many authors have already studied the case when the wave propagation speeds are equal to one another. \\
Note that, with our notation, the equality of the three wave propagation speeds is written as:

\begin{equation}
a = b = 1.
\label{eqwavespeeds}
\end{equation}

\noindent In the following, this case, which is not physically relevant, is excluded. We always assume the following condition: 

\begin{equation}
(H): (a,b) \neq (1,1) \; \mbox{and the cases} \; a=1, b \neq 1 ; \; a \neq 1, b = 1 ; \; a=b, a \neq 1 \; \mbox{are also excluded.}
\label{cond}
\end{equation}

\begin{remark}(Absence of physical meaning) \\
The assumption $a \neq b$ does not make sense physically since $a$ and $b$ coincide: they are both equal to $E/(\kappa' G)$.  \\
The method developed by Han and Xu in \cite{HanXu}, which is used in the following, requires this hypothesis (see Remarks \ref{rkequalvelocities} and \ref{rkHanXu}). But it is not clear why it would hold in the more general case treated by Han and Xu either. Lemma \ref{lemD} is the crucial point in their problem and ours. They should make assumptions which are not easy to formulate and which would also lead to physically irrelevant situations. \\
The authors are conscious of the fact that all what follows is only interesting from a mathematical point of view. 
\end{remark}

\noindent Indeed, the two moduli $E$ and $G$ are related by the expression 

$$G=\dfrac{E}{2(1+\nu)}$$

\noindent where $\nu\in(0,\frac{1}{2})$ is the Poisson ratio. Thereby, the study of the stabilisation when the speeds of propagation are equal to one another is considered only as mathematical study (see \cite{LiuRao}). That is why it is excluded in this paper. \\
\\
\noindent In the paper by J. E. Mu�oz Rivera and M.G. Naso (\cite{RN}), the authors consider the same system as us but with the boundary conditions $(BC_1)$:

\begin{equation} \label{BC1}
\left\{\begin{array}{lll}
\varphi(t,0) = \psi(t,1)= \omega(t,1)=0\\
(\varphi_x + \psi + l \omega)(t,1) = 0, \; (\omega_x-l \varphi)(t,0) = 0 \\ 
b c^{-1} \psi_x(t,0)= \gamma \psi_t(t,0)
\end{array}\right.
\end{equation}

\noindent where $\gamma$ is a strictly positive constant. (The problem they considered has been translated with our notation, so that the reader can see the analogies.) \\
They only have one pointwise dissipative mechanism (and not two) but they assume that $(a,b)=(1,1)$ i.e. that the three wave propagation speeds are equal to one another. They use multiplier methods to get exponential stability assuming a condition similar to (\ref{condnorm}) below. \\
They also show the same result with the boundary conditions $(BC_2)$ and $(a,b)=(1,1)$:

\begin{equation} \label{BC2}
\left\{\begin{array}{lll}
\varphi(t,0) =  \psi(t,1)= \omega(t,1)=0\\
(\varphi_x + \psi + l \omega)(t,1) = 0, \; \psi_x(t,0) = 0 \\ 
(\omega_x-l \varphi)(t,0) = 0 = \gamma \omega_t(t,0).
\end{array}\right.
\end{equation}

\noindent Earlier, in the paper by M.S. Alves, J.E.  Mu\~{n}oz Rivera, O.P. Vera Villagran and A. Rambaud (\cite{AVRR}), the authors consider the same system as us but with three boundary dissipation conditions:

\begin{equation} \label{BC3}
\left\{\begin{array}{llll}
\varphi(t,1)=\psi(t,1)=\omega(t,1)=0 \\
(\varphi_x(t,0) + \psi_x(t,0) + l \omega_x(t,0)= \gamma_1 \varphi_t(t,0) \\
b c^{-1} \psi_x(t,0)= \gamma_2 \psi_t(t,0) \\
a(\omega_x - l \varphi)(t,0)= \gamma_3 \omega_t(t,0).
\end{array}\right.
\end{equation}

\noindent In this problem, $\gamma_1$, $\gamma_2$ and $\gamma_3$ are strictly positive constants. The authors prove the exponential stability using the semi-group theory for linear operators and a result obtained by Pr\"uss. They have no assumptions on the wave propagation velocities ($a$ and $b$ with our notation) nor on $l$ (see our hypothesis (\ref{condnorm}) further, which is necessary to get a norm here, since we only put two dissipations).\\
\\
%%%%%%%%%%%%%%%%%%%%% Partie Biblio %%%%%%%%%%%%%%%%%%%%%%%%%%%%%%%%%%%%%%%%%%%%%%%%%%%%%%%%%%%
%% Avec les dampings dans les �quations
\noindent The Bresse system was also studied with frictional dissipative terms acting in all the equations. In \cite{Santos} e.g., the exponential decay of the solution is proved by using a method developed by Z. Liu and S. Zheng and their collaborators previously. The numerical computations are made by using the finite difference method to corroborate the theoretical results. In particular, the finite difference method in their case is locking free. \\
\\
More recently, in \cite{AkilBad}, the stabilization of a linear Bresse system with one singular local frictional damping acting in the longitudinal displacement, under fully Dirichlet boundary conditions, is investigated. Strong stability is first proved. Then, using a frequency domain approach combined with the multiplier method, the exponential stability of the solution is established if the three waves have the same speed of propagation. On the contrary, the energy of the system is proved to decay polynomially with rates $t^{-1}$ or $t^{-1/2}$. \\
\\
In \cite{AkilBadNicWeh}, the authors investigate the stabilization of a linear Bresse system with one discontinuous local internal viscoelastic damping of Kelvin-Voigt type acting on the axial force, under fully Dirichlet boundary conditions. First, using a general criteria of Arendt-Batty, they prove the strong stability of their system. Finally, using a frequency domain approach combined with the multiplier method, they prove that the energy of their system decays polynomially with different rates. \\
\\
\noindent An interesting situation is also that of the thermoelastic Bresse system which describes the motion of a linear planar, shearable thermoelastic beam. See \cite{LiuRao} for example. Z. Liu and B. Rao study the energy decay rate for this system. If the longitudinal motion and heat transfer are neglected, their model reduces to the well-known thermoelastic Timoshenko beam equations. The system they consider consists of three wave equations and two heat equations coupled in a certain pattern. The two wave equations about the longitudinal displacement and shear angle displacement are effectively damped by the dissipation from the two heat equations. Actually, the corresponding energy decays exponentially like the classical one-dimensional thermoelastic system. However, the third wave equation about the vertical displacement is only weakly damped. Thus the decay rate of the energy of the overall system was still unknown at the time. Z. Liu and B. Rao show that the exponentially decay rate is preserved when the wave speed of the vertical displacement coincides with the wave speed of longitudinal displacement or of the shear angle displacement. Otherwise, only a polynomial type decay rate can be obtained. These results are proved by verifying the frequency domain conditions.\\
\\
%%%  Papiers de Wael


%%%
\noindent The Bresse system is a valid model for arched beams, which reduces to the classical Timoshenko system, when the arch curvature $l$ is zero. Many papers are concerned with the boundary stabilization of the Timoshenko beam. \\
J. U. Kim and Y. Renardy proved the exponential stability of the system under two boundary controls in \cite{KR}. In \cite{Ammar2}, F. Ammar-Khodja, S. Kerbal and A. Soufyane studied the decay rate of the energy of the nonuniform Timoshenko beam with two boundary controls acting in the rotation-angle equation. Under the equal speed wave propagation condition (that is to say, if $l=0$ and $b=1$, in our modelization), they established exponential decay results up to an unknown finite dimensional space of initial data. In addition, they showed that the equal speed wave propagation condition is necessary for the exponential stability. \\
In \cite{AkilChiGhaWeh}, the indirect boundary stability and exact controllability of a one-dimensional Timoshenko system are investigated. In the first part of the paper, the Timoshenko system with only one boundary fractional damping is considered. The system is proved to be strongly stable but not uniformly stable. Using frequency domain arguments combined with the multiplier method, the authors prove that the energy polynomial decay rate depends on coefficients appearing in the PDE's and on the order of the fractional damping. Moreover, under the equal speed propagation condition, they obtain the optimal polynomial energy decay rate. In the second part of the paper, the indirect boundary exact controllability of the Timoshenko system with mixed Dirichlet-Neumann boundary conditions and boundary control is studied. Using non-harmonic analysis, the authors first establish a weak observability inequality, which depends on the ratio of the wave propagation speeds. Next, using the HUM method, they prove that the system is exactly controllable in appropriate spaces and that the control time can be small. \\
\\
In \cite{mr1} and in \cite{EECT18}, the same system as the one considered in the following was studied but with different boundary conditions: $u(0,t)=0$ in the first one and $u_x(0,t)=0$ in the second one. In \cite{mr1}, under the equal speed condition ($b=1$ and $l=0$ in our modelization) and if $c$ is outside a discrete set of exceptional values, a polynomial energy decay rate is obtained using a spectral analysis. On the other hand, if $\sqrt{b}$ is a rational number and if $c$ is outside another discrete set of exceptional values, a polynomial type decay rate is proved to hold using a frequency domain approach. \\
The main goal of \cite{EECT18} is to obtain the energy decay rate if $\sqrt{b}$ is a rational number but without assumptions on $c$ (except $(C_1)$ and/or $(C_2)$, introduced in Theorem 2.2, which are necessary for strong stability) and with the boundary condition $u_x(0,t)=0$ instead of $u(0,t)=0$. The result is proved by means of suitable estimates of the resolvent operator norm along the imaginary axis. The technique which is chosen involves the Laplace transform which is an innovative technical choice. \\ 
If $\sqrt{b}$ is not a rational number, the obtention of the energy decay rate is reduced to a non trivial number theory problem and it is conjectured that the decay rate of the energy in that case is very small. The problem is still open. \\
\\
\noindent The same method could theoretically be used to study System (\ref{bresse}) with conditions (\ref{IC}) and (\ref{BC}). In practice, since the calculations were mainly explicit in \cite{EECT18}, it is far too complicated. The polynomial $\Delta$ defined by $(19)$ is a polynomial of degree two, whose roots are explicitly known as functions of the variable $\mu$. For the Bresse system, this polynomial would be of degree three. The asymptotic behaviour of the functions $J_k$ for large values of $\mu$, is established using a rewriting of $\Phi$ in Lemma 4.1 as well as explicit decompositions of rational fractions. The asymptotic behaviour of $q_1$ and $q_2$ is explicitly computed in the proof of Theorem 5.1. All this was possible since the involved expressions were not so awful. \\
That is why we proceed in a different way here for the Bresse system. \\
\\
%%%  Plan du papier
\noindent The paper is structured as follows: the existence, uniqueness and strong stability of the solution of System (\ref{bresse}) with the initial conditions (\ref{IC}) and the boundary conditions (\ref{BC}) are established in the first section, using the maximal dissipativeness of the underlying operator. \\
Then, in Section $3$, the asymptotic behaviour of the eigenvalues of this operator is computed, using the same method as used by Z-J. Han and G-Q Xu in \cite{HanXu}. Note that this method involves the matrix $M_1$ defined by (\ref{decompoM}) and it requires that the eigenvalues of this matrix are all simple. In our situation, it is the case when the wave propagation velocities are distinct, which is the case we are interested in. If some of the wave propagation velocities coincide, the multiplicity of some of the eigenvalues of the matrix $M_1$ may be equal to $2$ or $3$. Z-J. Han and G-Q Xu have neglected this problem in \cite{HanXu}. \\
The asymptotic behaviour of the eigenfunctions of the operator is given in Section $4$.
%which allows to prove that they form a Riesz basis of the state Hilbert space. 
Since the imaginary axis is an asymptote for one family of eigenvalues, the decreasing of the energy can not be exponential. A lot more work is required to get the decay rate of the energy. \\
\\
\noindent The main result of the paper is summed up in the following Theorem:

\bthe \label{theig}(Asymptotic behaviour of a family of eigenvalues of the underlying operator)   \\
Assume Condition (\ref{cond}) on $a$ and $b$ as well as Condition (\ref{condnorm}). \\
The asymptotic behaviour of one family of eigenvalues of the operator $\caA$ defined by (\ref{defDA}) and (\ref{defA}) is given by:

$$\mu_n = i n \pi + O(n^{-1}), n \in \Z$$

\noindent i.e. the imaginary axis is an asymptote for these eigenvalues.
\ethe
 

%%%%%%%%%%%%%%%%%%%%%%%%%%%%%%%%%%%%%%%%%%%%%%%%%%%%%%%%%%%%%%%%%%%%%%%%%%%%%%%%%%%%%%%%%%%%%%%%%%%%%%%%%%%%%%%%%%%%%%%%%%%%%%%%%%%%%%%%%%%%%%
%%%%%%%%%%%%%%%%%%%%%%%%%%%%%%%%%%%%%%%%%%%%%%%%%%%%%%%%%%%%%%%%%%%%%%%%%%%%%%%%%%%%%%%%%%%%%%%%%%%%%%%%%%%%%%%%%%%%%%%%%%%%%%%%%%%%
\section{Well-posedness and strong stability} \label{sec2}
%\setcounter{equation}{0}

\noindent In this section we study the existence, uniqueness and strong stability of the solution of System (\ref{bresse}) with the initial conditions (\ref{IC}) and the boundary conditions (\ref{BC}).\\ 
\\
\noindent The energy space $\mathcal{H}$ is defined as follows

\be \label{defHRL}
H_0^1(0;1) = \{ \varphi \in H^1(0;1) / \varphi(0) = \varphi(1)=0 \}, H_L^1(0;1) = \{ \psi \in H^1(0;1) / \psi(0)=0 \}
\ee

\be \label{defH}
\mathcal{H} = H_0^1(0;1) \times L^2(0;1) \times H_L^1(0;1) \times L^2(0;1) \times H_L^1(0;1) \times L^2(0;1)
\ee

\noindent with the inner product defined by

\beq
(U_1,U_2)_{\caH} = \begin{split}
& \displaystyle\int_0^1 (\eta_1 \overline{\eta_2} + c^{-1} \theta_1 \overline{\theta_2} + \zeta_1 \overline{\zeta_2} + b c^{-1} \psi_{1,x} \overline{\psi_{2,x}} + (\varphi_{1,x} + \psi_1 + l \omega_1) (\overline{\varphi_{2,x}} + \overline{\psi_2} + l \overline{\omega_2}) \\
& + a (\omega_{1,x} - l \varphi_1)(\overline{\omega_{2,x}} - l \overline{\varphi_2}))dx,
\end{split}
\label{innprod}
\eeq

\noindent for all $U_1 = (\varphi_1,\eta_1, \psi_1, \theta_1, \omega_1, \zeta_1)$, $U_2 = (\varphi_2,\eta_2, \psi_2, \theta_2, \omega_2, \zeta_2)$ $\in$ $\caH$. \\
\\
The functional defined by (\ref{innprod}) does not always define a norm. Since the proof is the same as in the paper by J. E. Mu�oz Rivera and M.G. Naso (\cite{RN}), we do not give all the details of the proof.

%%%%%%%%%%%%%%%%%%%%%%%         Lemma norm            %%%%%%%%%%%%%%%%%%%%%%%%%%%%%%%%%
\blem
Assume the following condition:

\beq
(2 k + 1) \dfrac{\pi}{2} \neq l, \forall k \in \N.
\label{condnorm}
\eeq

The functional $(U,U)_{\caH}^{\frac12}$ induced by (\ref{innprod}) defines a norm which is equivalent to the usual norm of $\caH$.
\elem
%%%%%%%%%%%%%%%%%%%%%%%%%%%%%%%%%%%%%%%%%%%%%%%%%%%%%%%%%%%%%%%%%%%%%%%%%%%%%%%%%%%%%%%%%%%%%%%%

\proof 
\noindent If $(U,U)_{\caH} = 0$, then $\eta = 0$, $\theta = 0$, $\zeta = 0$ and

$$\psi_x=0, \varphi_{x} + \psi + l \omega = 0, \; \omega_{x} - l \varphi = 0.$$

\noindent Since $\psi \in H_L^1(0;1)$, the first equation implies $\psi =0$. \\
From the other two equations follows: $\omega_{xx} + l^2 \varphi = 0$. The solution is: $\omega(x) = K \sin(l x)$, with $K \in \R$, since $\omega(0)=0$. \\
Now $\varphi_{x} + \psi + l \omega = 0$ with $\psi = 0$ and $\omega(x) = K \sin(l x)$ leads to: $\varphi(x) = C_1 \cos(l x) + C_2$. Using the fact that $\varphi(0)=\varphi(1) = 0$, we get: either $(C_1,C_2)=(0,0)$ and $\varphi = \psi = \omega = 0$. Then $U= 0$ and the functional induced by (\ref{innprod}) defines a norm which is equivalent to the usual norm of $\caH$. \\
Or $(C_1,C_2) \neq (0,0)$ and $\cos(l)=1$. This can only hold when Condition $(\ref{condnorm})$ fails.  \fin \\
\\
\noindent \bf From now on, Condition (\ref{condnorm}) is supposed to hold. \rm \\
\\
\noindent For shortness we denote by $\|\,.\,\|$   the  $L^2(\Omega)$-norm. \\
\\
\noindent Now, we define a linear unbounded operator $\mathcal{A}:D(\mathcal{A})\rightarrow\caH$ by:

\begin{eqnarray} \label{defDA}
\nonumber D(\mathcal{A}) = \bigg \{ U = (\varphi,\eta, \psi, \theta, \omega, \zeta) \in \mathcal{H} : (\varphi,\psi,\omega) \in H^2(0,1) \times H^2(0;1) \times H^2(0;1),\, \eta \in H_0^1(0;1), \theta \in H_L^1(0;1) \\
\zeta \in H_L^1(0;1), b c^{-1} \psi_{x}(1) = - \beta_1 \theta(1), a(\omega_x(1) - l \varphi(1)) = - \beta_2 \zeta(1)   \bigg \} 
\end{eqnarray}


\begin{eqnarray} \label{defA}
\nonumber \mathcal{A}(\varphi,\eta, \psi, \theta, \omega, \zeta) = \big(\eta, (\varphi_x + \psi + l\omega)_x + al (\omega_x - l \varphi), \theta, b \psi_{xx} - c(\varphi_x + \psi + l\omega), \\
\zeta, a (\omega_x - l \varphi)_x - l(\varphi_x + \psi + l\omega) \big), \forall\, U = (\varphi,\eta, \psi, \theta, \omega, \zeta) \in D(\mathcal{A}).
\eeq


\noindent Then we rewrite formally System (\ref{bresse}) into the evolution equation

\begin{equation}
\left\{\begin{array}{ll}
U_t = \mathcal{A}U,\\
U(0) = U_0, \ \ \ \ U_0\in \caH
\label{abstr}
\end{array}\right.
\end{equation}
with  $U = (u,u_t)$.

\begin{proposition}
The operator $\caA$ is m-dissipative in the energy space $\caH$.
\end{proposition}

%%%%%%%%%%%%%%%%%%%%%%%%%%%%%%%%%%%%%%%%%%%%%%%%%%%%%%%%%%%%%%%%%%%%%%%%%%%%%%%%%%%%%%%%%%%

\proof On the first hand, the dissipativeness holds since we can check using integrations by parts: 

\begin{equation}
\Re (\caA U, U)_{\caH} = -\beta_1 \vert\psi_t(t,1)\vert^2 - \beta_2 \vert \omega_t(t,1) \vert^2 \leq 0,  \forall \, U = (\varphi,\varphi_t, \psi, \psi_t, \omega, \omega_t) \in \caD(\caA). \label{AUU}
\end{equation} 

\noindent On the other end, the maximality is proved in the following way (cf. the paper by M. Bassam, D. Mercier, S. Nicaise and A. Wehbe: \cite{mr1}).\\
\\
\noindent Let $f := (f_1;f_2;f_3;f_4;f_5,f_6) \in \caH$. We look for $U:=(\varphi,\eta, \psi, \theta, \omega, \zeta)  \in \mathcal{D}(\mathcal{A})$ solution of 

\begin{equation}\label{eqmaxmon}
- \mathcal{A} U =  f
\end{equation} 

\noindent or equivalently

\begin{equation} \label{f}
\left\{\begin{array}{llllll}
f_1 = - \eta \\
f_2 = - (\varphi_x + \psi + l \omega)_x  - a l (\omega_x - l \varphi) \\
f_3 = - \theta \\
f_4 = - b \psi_{xx} + c (\varphi_x + \psi + l \omega) \\
f_5 = - \zeta \\
f_6 = - a (\omega_x - l \varphi)_x + l (\varphi_x + \psi + l \omega)  
\end{array} \right. 
\end{equation}

\noindent Assume that such a solution $U$ exists. We multiply the second identity by the product $c \overline{\varphi_1}$ with $\varphi_1 \in H_R^1(0;1)$, the fourth identity by the function $\overline{\psi_1}$ with $\psi_1 \in H_L^1(0;1)$ and the sixth identity by the product $c \overline{\omega_1}$ with $\omega_1 \in H_L^1(0;1)$. Then, we take the sum of these three identities and integrate over $(0;1)$. At last, three integrations by parts, taking into account the boundary conditions ($U \in \mathcal{D}(\mathcal{A})$), like in \cite{mr1}, lead to: 

\begin{eqnarray} \label{maxeq}
\nonumber \int_0^1 \left[ b \psi_{x} \overline{\psi_{1,x}} + c (\varphi_x + \psi + l \omega)(\overline{\varphi_{1,x}} + \overline{\psi_1} + l \overline{\omega_1}) + a c (\omega_x - l \varphi)(\overline{\omega_{1,x}} - l \overline{\varphi_1}) \right](x) dx = \\
\int_0^1 \left[ c f_2 \overline{\varphi_1} + f_4 \overline{\psi_1} +  c f_2 \overline{\omega_1} \right] (x) dx + c \beta_1 f_3(1) \overline{\psi_1}(1) + c \beta_2 f_5(1) \overline{\omega_1}(1).
\end{eqnarray} 

\noindent This problem has a unique solution $U \in V := H_0^1(0;1) \times H_L^1(0;1)) \times H_L^1(0;1)$ by Lax-Milgram's lemma, because the left-hand side  of (\ref{maxeq}) defines a coercive sesquilinear form on $V$.\\
\\
\noindent By taking test functions of the form $(u, 0, 0)$, $(0, y, 0)$ and $(0,0,z)$ with $u$, $y$ and $z$ in $\mathcal{D}(0,1)$, we deduce that:

\begin{equation}
\left\{\begin{array}{llllll}
\varphi_{xx} = - f_2 -(\psi + l \omega)_x - a l (\omega_x - l \varphi) \; \mbox{in the distributional sense} \\
\psi_{xx} = - b^{-1} f_4 + b^{-1} c (\varphi_x + \psi + l \omega) \; \mbox{in the distributional sense}\\
\omega_{xx} = - a^{-1} f_6  + l \varphi_x + a^{-1} l (\varphi_x + \psi + l \omega)  \; \mbox{in the distributional sense} 
\end{array} \right. 
\end{equation}

\noindent This directly implies that  $\varphi$, $\psi$ and $\omega$ belong to $H^2(0;1)$ since $f_2$, $f_4$ and $f_6$ are in $L^2(0;1)$. \\
\\
\noindent Now, we choose $\overline{\varphi_1}=0$ and $\overline{\omega_1} = 0$ in the variational formulation (\ref{maxeq}), integrate by parts and use the fourth equation of (\ref{f}) as well as the fact that $\overline{\psi_1}(0)=0$ to get:

\begin{eqnarray} \label{maxeq2}
\nonumber \int_0^1 \left[ b \psi_{x} \overline{\psi_{1,x}} + c (\varphi_x + \psi + l \omega)\overline{\psi_1} \right](x) dx = \int_0^1 f_4 \overline{\psi_1}(x) dx + c \beta_1 f_3(1) \overline{\psi_1}(1)  \\
\Leftrightarrow - \int_0^1 \left[ b \psi_{xx} \overline{\psi_1} + c (\varphi_x + \psi + l \omega)\overline{\psi_1} \right](x) dx - b [ \psi_{x} \overline{\psi_1}]_{0}^{1} = \int_0^1 f_4 \overline{\psi_1}(x) dx + c \beta_1 f_3(1) \overline{\psi_1}(1)  \\
\Leftrightarrow c \beta_1 f_3(1) \overline{\psi_1}(1) = b \psi_x(1) \overline{\psi_1}(1).
\end{eqnarray} 

\noindent Since this must hold for any function $\psi_1$ in $H_L^1(0;1)$, it follows:

$$b \psi_x(1) = c \beta_1 f_3(1) = - c \beta_1 \theta(1) \Leftrightarrow b c^{-1} \psi_x(1) =  - \beta_1 \theta(1).$$

\noindent Then, we choose $\overline{\psi_1}=0$ in the variational formulation (\ref{maxeq}), integrate by parts and use the second and sixth equations of (\ref{f}) as well as the fact that $\overline{\varphi_1}(0)= \overline{\varphi_1}(1)=  \overline{\omega_1}(1) = 0$ to get, with the same arguments as before:

$$a(\omega_x(1) - l \varphi(1)) = \beta_2 f_5(1) = - \beta_2 \zeta(1).$$

\noindent Consequently, the unique solution $U \in V := H_0^1(0;1) \times H_L^1(0;1)) \times H_L^1(0;1)$ satisfying (\ref{maxeq}) which we had obtained earlier, is in fact in $\mathcal{D}(\mathcal{A})$. This finishes the proof of maximality.             \fin

\medskip


%%%%%%%%%%%%%%%%%%%%%%%%%%%%%%%%%%%%%%%%%%%%%%%%%%%%%%%%%%%%%%%%%%%%%%%%%%%%%%%%%%%%%%%%%


\noindent Using Lumer-Phillips Theorem (see the book by A. Pazy, \cite{P}, Theorem 1.4.3), the operator $\caA$ generates a $C_0$-semigroup of contractions $e^{t\caA}$ on
$\caH$. Then, we have the following result.

\bthe(Existence and uniqueness)\\
(1) If $U_0$ $\in$ $\caD(\caA)$, then   system (\ref {abstr}) has a unique strong solution
$$U  \in C^0(\R_+,\caD(\caA))\cap C^1(\R_+,\caH).$$
(2) If $U_0$ $\in$ $\caH$, then   system (\ref {abstr}) has a unique weak solution\\
$$U \in C^0(\R_+,\caH).$$
\ethe

%%%%%%%%%%%%%%%%%%%%%%%%%%%%%%%%%%%%%%%%%%%%%%%%%%%%%%%%%%%%%%%%%%%%%%%%%%%%%%%%%%%%%
\noindent Now, we have the following general strong stability result. 

\bthe \label{thstrong}(Strong stability)   \\
System (\ref{bresse})-(\ref{IC})-(\ref{BC}) is strongly stable, i.e. for any solution $U$ of (\ref{abstr}) with initial data $U_0 \in\caH$, it holds 
$$\lim_{t\rightarrow\infty }E(t)=0,$$ where $E(t)$ is defined by (\ref{def.energy}). 
\ethe

%%%%%%%%%%%%%%%%%%%%%%%%%%%%%%%%%%%%%%%%%%%%%%%%%%%%%%%%%%%%%%%%%%%%%%

\proof Since $\caA$ generates a contraction semigroup and its resolvent is compact in $\caH$, using Arendt-Batty Theorem (see \cite{AB}, p. 837), system (\ref{bresse})-(\ref{IC})-(\ref{BC}) is strongly stable if and only if $\caA$ does not have purely imaginary eigenvalues. \\
\\
\noindent Assume that $\caA$ has a purely imaginary eigenvalue denoted by $i \mu$ with $\mu \in \R$. Then there exists $U:=(\varphi,\eta, \psi, \theta, \omega, \zeta) \neq (0;0;0;0;0;0) \in \mathcal{D}(\mathcal{A})$ such that $\caA U = i \mu U$. Using (\ref{AUU}) leads to:   

\begin{equation}
\Re \left[ i \mu (U, U)_{\caH} \right] = 0 = -\beta_1 \vert \theta(1) \vert^2 - \beta_2 \vert \zeta(1) \vert^2. \label{strstab}
\end{equation} 

\noindent This implies $|\theta(1)| = 0$ and $|\zeta(1)| = 0$ (since $\beta_1>0$ and $\beta_2 > 0$). Now $\theta$ and $\zeta$ are continuous on $(0;1)$ and they vanish at $0$ since $U$ belongs to $D(\caA)$. Thus $\theta \equiv \zeta \equiv 0$ and since $\theta = i \mu \psi$ and $\zeta = i \mu \omega$, $\psi \equiv 0$ and $\omega \equiv 0$. \\
Since $\caA U = i \mu U$ can be rewritten as the following system:

\begin{equation} \label{AUU2}
\left\{\begin{array}{llllll}
i \mu \varphi = \eta \\
i \mu \eta = (\varphi_x + \psi + l \omega)_x  + a l (\omega_x - l \varphi) \\
i \mu \psi = \theta \\
i \mu \theta = b \psi_{xx} - c (\varphi_x + \psi + l \omega) \\
i \mu \omega = \zeta \\
i \mu \zeta = a (\omega_x - l \varphi)_x - l (\varphi_x + \psi + l \omega)  
\end{array} \right. 
\end{equation}

\noindent Using the fourth or the sixth equation of (\ref{AUU2}) leads to $\varphi_x \equiv 0$ or equivalently to: $\varphi$ is a constant function. Now $\varphi$ is continuous on $(0;1)$. Since $U$ belongs to $D(\caA)$, $\eta \in H_0^1(0;1)$ and, since $i \mu \varphi = \eta$, $\varphi$ vanishes at $0$ and $1$. All that requires $\varphi$ to vanish and so $\eta$. \\
To conclude, $\theta \equiv 0$, $\zeta \equiv 0$, $\psi \equiv 0$, $\omega \equiv 0$, $\varphi \equiv 0$ and $\eta \equiv 0$. \\
This contradicts the fact that $U \neq (0;0)$.   \fin

%%%%%%%%%%%%%%%%%%%%%%%%%%%%%%%%%%%%%%%%%%%%%%%%%%%%%%%%%%%%%%%%%%%%%%%%%%%%%%%%%%%%%%%%%%%%%%
%%%%%%%%%%%%%%%%%%%%%%%%%%%%%%%%%%%%%%%%%%%%%%%%%%%%%%%%%%%%%%%%%%%%%%%%%%%%%%%%%%%%%%%%%%%%%%%%%%%%%%%%%%%%%%%%%%%%%%%%%%%%%%%%%%%%%%%%%%%%%%
\section{Eigenvalues of the operator $\caA$} \label{sec3}
%%%%%%%%%%%%%%%%%%%%%%%%%%%%%%%%%%%%%%%%%%%%%%%%%%%%%%%%%%%%%%%%%%%%%%%%%

\noindent Before we prove Theorem \ref{theig} which gives the eigenvalues of the operator $\caA$ as solutions of the characteristic equation (\ref{chareq0}) and the asymptotic behaviour of some of them, we need some technical lemmas. 


%%%%%%%%%%%%%%%%%%%%%%%     Lemma matrices     %%%%%%%%%%%%%%%%%%%%%%%%%%%%   
\begin{lemma} (Product with a diagonal matrix) \\ \label{lemD}
Let $A$ and $D$ be square matrices of order $n$, with $n \geq 2$. Assume that $D$ is diagonal.

\begin{itemize}
\item The diagonal terms of $A \cdot D - D \cdot A$ vanish. 
\item If the terms of $D$ are all different i.e. $d_{ii} \neq d_{kk}, \forall (i,k) \in [1,n]^2$, then 
$$A \cdot D = D \cdot A \implies \; \mbox{$A$ is diagonal}.$$
\end{itemize}

\end{lemma}

\proof
\noindent The proof is left to the reader. \fin \\
\\
\bthe \label{theig}(Characterization and asymptotic behaviour of the eigenvalues)   \\
Assume Condition (\ref{condnorm}). Define by $D$ the diagonal matrix $D:=diag \left( 1,-1,\sqrt{\frac{1}{a}}, -\sqrt{\frac{1}{a}}, \sqrt{\frac{1}{b}}, -\sqrt{\frac{1}{b}} \right)$ and by $S(\mu)$ the square matrix of order $6$ given by:

\begin{equation} \label{Sdef1}
S(\mu):= \sum_{k \geq 0} \dfrac{S_{-k}(1)}{\mu^k} e^{\mu D}
\end{equation}

\noindent where the terms of $S_{-k}(1)$ are computable via an iterative procedure explained in the proof. \\
The eigenvalues of the operator $\caA$ defined by (\ref{defDA}) and (\ref{defA}) satisfy the characteristic equation:

\begin{equation} \label{chareq0}
\alpha_6 (\alpha_6 \gamma_7 - \gamma_6 \alpha_7) \tau = 0.
\end{equation}

\noindent where $\alpha_6, \alpha_7, \alpha_8, \gamma_6, \gamma_7, \gamma_8, \tau$ are all functions of $\mu$ defined by:

\begin{equation} \label{def_alpha_gamma0} 
\left \{ \begin{array}{lllllll} 
\alpha_6:= -(\beta_1 c + \sqrt{b}) \sqrt{b} (s_{51} + s_{52}) - (\beta_1 c - \sqrt{b}) \sqrt{b} (s_{61} + s_{62}) \\
\alpha_7:= -(\beta_1 c + \sqrt{b}) \sqrt{b} (s_{55} - s_{56}) - (\beta_1 c - \sqrt{b}) \sqrt{b} (s_{65} - s_{66}) \\
\alpha_8:= -(\beta_1 c + \sqrt{b}) \sqrt{b} (s_{53} - s_{54}) - (\beta_1 c - \sqrt{b}) \sqrt{b} (s_{63} - s_{64}) \\
\gamma_6:= -(\beta_2 + \sqrt{a}) \sqrt{a} (s_{31} + s_{32}) - (\beta_2 - \sqrt{a}) \sqrt{b} (s_{41} + s_{42}) \\
\gamma_7:= -(\beta_2 + \sqrt{a}) \sqrt{a} (s_{35} - s_{36}) - (\beta_2 - \sqrt{a}) \sqrt{b} (s_{45} - s_{46}) \\
\gamma_8:= -(\beta_2 + \sqrt{a}) \sqrt{a} (s_{33} - s_{34}) - (\beta_2 - \sqrt{a}) \sqrt{b} (s_{43} - s_{44})  
\end{array} \right.
\end{equation}

\begin{equation} \label{def_tau0} 
\tau:= \det \left( \begin{array}{cc} 
\alpha_6 \gamma_7 - \gamma_6 \alpha_7 & \alpha_6 \gamma_7 - \gamma_6 \alpha_7 \\
\alpha_6 (h_{27} - h_{17}) - \alpha_7 (h_{26} - h_{16}) & \alpha_6 (h_{28} - h_{18}) - \alpha_8 (h_{26} - h_{16})
\end{array} \right)
\end{equation}

\noindent with $h_{16}:= s_{11} + s_{12}$, $h_{26}:= s_{21} + s_{22}$, $h_{17}:= s_{15} - s_{16}$, $h_{27}:= s_{25} - s_{26}$ and $s_{ij}:= s_{ij}(\mu) = (S(\mu))_{ij}$. \\
\\
A first family of eigenvalues satisfies $\alpha_6(\mu)=0$. Assume Condition (\ref{cond}) on $a$ and $b$, then the asymptotic behaviour of this family is given by $\mu_n = i n \pi + O(n^{-1})$, $n \in \Z$ i.e. the imaginary axis is an asymptote for these eigenvalues.
\ethe

\proof First of all, the operator $\caA$ has already been proved to have no imaginary eigenvalues (cf. the proof of Theorem \ref{thstrong}). In particular, $0$ is not an eigenvalue. \\
\\
\noindent The computation of the eigenelements of $\caA$ requires the search for a solution $U =(\varphi, \eta, \psi, \theta, \omega, \zeta) \neq 0$ of $\caA U = \mu U$ ($\mu \in \C^*$), which can be rewritten as the following system:

\begin{equation} \label{eigenpb}
\left\{\begin{array}{llllll}
\mu \varphi = \eta \\
\mu \eta = (\varphi_x + \psi + l \omega)_x  + a l (\omega_x - l \varphi) \\
\mu \psi = \theta \\
\mu \theta = b \psi_{xx} - c (\varphi_x + \psi + l \omega) \\
\mu \omega = \zeta \\
\mu \zeta = a (\omega_x - l \varphi)_x - l (\varphi_x + \psi + l \omega)  
\end{array} \right. 
\end{equation}

\noindent It is equivalent to:

\begin{equation} \label{eigenpb2}
\left\{\begin{array}{llllll}
- \varphi_{xx} - \psi_x - (a + 1) l \omega_x + (\mu^2 + a l^2) \varphi = 0   \\
- b \psi_{xx} + c \varphi_x + c l \omega + (\mu^2 + c) \psi = 0 \\
-a \omega_{xx} + (a + 1) l \varphi_x + l \psi + (\mu^2 + l^2) \omega = 0 \\
\eta = \mu \varphi \\
\theta = \mu \psi \\
\zeta = \mu \omega  \\  
\end{array} \right. 
\end{equation}

\noindent Since $\mu \neq 0$, we denote by $W$ the vector 

\begin{equation} \label{defW}
W(\mu,x) : = \left(\varphi, \psi, \omega, \dfrac{\varphi_x}{\mu}, \dfrac{\psi_x}{\mu},\dfrac{\omega_x}{\mu} \right)^t.
\end{equation}

\noindent The system formed by the first three equations of (\ref{eigenpb2}) is written in the matricial form 

\begin{equation} \label{Weq}
W_x(\mu,x) = M(\mu) W(\mu,x)
\end{equation}

\noindent with $M(\mu)$ defined as:

$$M(\mu): = \left( \begin{array}{cccccc} 
0 & 0 & 0 & \mu & 0 & 0 \\
0 & 0 & 0 &  0 & \mu & 0 \\
0 & 0 & 0 & 0 & 0 & \mu \\
\mu + \frac{a l^2}{\mu} & 0 & 0 & 0 & -1 & -(a+1)l \\
0 & \frac{\mu}{b} + \frac{c}{\mu b} & \frac{cl}{\mu b} & \frac{c}{b} & 0 & 0 \\
0 & \frac{l}{\mu a} & \frac{\mu}{a} + \frac{l^2}{\mu a} & \frac{(a+1)l}{a} & 0 & 0 
\end{array} \right).$$

\noindent The idea is to proceed like Z-J. Han and G-Q Xu in Step $1$ of Section $4$ of \cite{HanXu}. The matrix $M(\mu)$ is decomposed as: 

\begin{equation} \label{decompoM}
M(\mu)= \mu M_1 + M_0 + \frac{1}{\mu} M_{-1}
\end{equation}

\noindent where $M_1$, $M_0$ and $M_{-1}$ are independent matrices of $\mu$ defined by: 

\begin{equation} \label{defM1}
M_1: = \left( \begin{array}{cccccc} 
0 & 0 & 0 & 1 & 0 & 0 \\
0 & 0 & 0 &  0 & 1 & 0 \\
0 & 0 & 0 & 0 & 0 & 1 \\
1 & 0 & 0 & 0 & 0 & 0 \\
0 & \frac{1}{b} & 0 & 0 & 0 & 0 \\
0 & 0 & \frac{1}{a} & 0 & 0 & 0 
\end{array} \right), M_0: = \left( \begin{array}{cccccc} 
0 & 0 & 0 & 0 & 0 & 0 \\
0 & 0 & 0 &  0 & 0 & 0 \\
0 & 0 & 0 & 0 & 0 & 0 \\
0 & 0 & 0 & 0 & -1 & -(a+1)l \\
0 & 0 & 0 & \frac{c}{b} & 0 & 0 \\
0 & 0 & 0 & \frac{(a+1)l}{a} & 0 & 0 
\end{array} \right)
\end{equation}

\begin{equation} \label{defMmoins1}
\mbox{and} \; M_{-1}: = \left( \begin{array}{cccccc} 
0 & 0 & 0 & 0 & 0 & 0 \\
0 & 0 & 0 &  0 & 0 & 0 \\
0 & 0 & 0 & 0 & 0 & 0 \\
al^2 & 0 & 0 & 0 & 0 & 0 \\
0 & \frac{c}{b} & \frac{cl}{b} & 0 & 0 & 0 \\
0 & \frac{l}{a} & \frac{l^2}{a} & 0 & 0 & 0 
\end{array} \right).
\end{equation}

\noindent The matrix $M_1$ is diagonalized to simplify the computations. Defining $P$ and $D$ as:

\begin{equation} \label{PdefD}
P: = \left( \begin{array}{cccccc} 
1 & -1 & 0 & 0 & 0 & 0 \\
0 & 0 & 0 &  0 & 1 & 1 \\
0 & 0 & 1 & 1 & 0 & 0 \\
1 & 1 & 0 & 0 & 0 & 0 \\
0 & 0 & 0 & 0 & \sqrt{\frac{1}{b}} & -\sqrt{\frac{1}{b}} \\
0 & 0 & \sqrt{\frac{1}{a}} & -\sqrt{\frac{1}{a}} & 0 & 0 
\end{array} \right), \; D:= \left( \begin{array}{cccccc} 
1 & 0 & 0 & 0 & 0 & 0 \\
0 & -1 & 0 &  0 & 0 & 0 \\
0 & 0 & \sqrt{\frac{1}{a}} & 0 & 0 & 0 \\
0 & 0 & 0 & -\sqrt{\frac{1}{a}} & 0 & 0 \\
0 & 0 & 0 & 0 & \sqrt{\frac{1}{b}} & 0 \\
0 & 0 & 0 & 0 & 0 & -\sqrt{\frac{1}{b}}
\end{array} \right).
\end{equation}

\noindent a straightforward calculation leads to: $P^{-1} M_1 P = D$. \\
\\


\noindent The solution of (\ref{Weq}) is $W(\mu, x) = \exp[M(\mu) x] W(\mu,0)$. Setting 

\begin{equation} \label{defZ}
Z(\mu,x) =P^{-1} W(\mu, x)
\end{equation}

\noindent and using the decomposition $M(\mu)= \mu M_1 + M_0 + \frac{1}{\mu} M_{-1}$, it holds:

\begin{eqnarray} \label{ZtoS}
Z(\mu, x) = \exp \left[ x \left( \mu D + P^{-1} M_0 P + \dfrac{1}{\mu} P^{-1} M_{-1} P \right) \right] Z(\mu,0) \\
\Leftrightarrow Z(\mu, x) = \exp \left[ x \mu D \left(I + \mu^{-1} D^{-1} P^{-1} M_0 P + \mu^{-2} D^{-1} P^{-1} M_{-1} P \right) \right] Z(\mu,0)
\end{eqnarray} 

\noindent Denoting by $N_0$ the matrix $D^{-1} P^{-1} M_0 P$ and $N_{-1}$ the matrix $D^{-1} P^{-1} M_{-1} P$, it follows:

\begin{eqnarray} \label{ZtoS2}
\nonumber Z(\mu, x) = \sum_{i \geq 0} \dfrac{1}{i !} x^i \mu^i D^i \left( I + \mu^{-1} N_0 + \mu^{-2} N_{-1} \right)^i  Z(\mu,0) \\
\nonumber \Leftrightarrow Z(\mu, x) = \sum_{i \geq 0} \dfrac{1}{i !} x^i \mu^i D^i \left( I + \dfrac{\Sigma_1}{\mu} + \dfrac{\Sigma_2}{\mu^2} + \ldots + \dfrac{\Sigma_{2i}}{\mu^{2i}}\right) Z(\mu,0) \\
\Leftrightarrow Z(\mu, x) = \sum_{k \geq 0} \dfrac{S_{-k}(x)}{\mu^k} \exp(\mu D x) Z(\mu,0).
\end{eqnarray} 

\noindent where $S_{-k}(x)$ is a square matrix of order $6$ which can be computed step by step as we will see further in the proof. \\
\\
\noindent Note that (\ref{ZtoS2}) implies: $S_0(0) = I, \, S_{-k}(0)=0, \forall k \geq 0$. \\
\\
\noindent The characteristic equation will follow from the expression of the boundary conditions at $0$ and $1$. Taking into account $\varphi(0)= \varphi(1) = \psi(0) = \omega(0) =0$, the last equation of (\ref{ZtoS2}) with $x=1$ is written as:

\begin{equation} \label{sysBC}
P^{-1} \left(0, \psi(1), \omega(1), \dfrac{\varphi_x(1)}{\mu}, \dfrac{\psi_x(1)}{\mu},\dfrac{\omega_x(1)}{\mu} \right)^t = S(\mu) P^{-1} \left(0,0,0, \dfrac{\varphi_x(0)}{\mu}, \dfrac{\psi_x(0)}{\mu},\dfrac{\omega_x(0)}{\mu} \right)^t.
\end{equation} 

\noindent where 

\begin{equation} \label{Sdef}
S(\mu):= \sum_{k \geq 0} \dfrac{S_{-k}(1)}{\mu^k} e^{\mu D}
\end{equation}

\noindent is a square matrix of order $6$. Setting 

\begin{equation} \label{defX}
X:= \left(\varphi_x(1), \psi_x(1), \omega_x(1), \psi(1), \omega(1), \varphi_x(0), \psi_x(0), \omega_x(0) \right)^t.
\end{equation} 

\noindent and adding the second line of (\ref{BC}) (the dissipation conditions at $x=1$), we obtain a system which can be written in the matricial form 

\begin{equation} \label{syschareq}
H(\mu) X = 0
\end{equation} 

\noindent with $H(\mu)$ the square matrix of order $8$ defined by:

\begin{equation} \label{matrixH} 
H(\mu): = \left( \begin{array}{cccccccc} 
-1 & 0 & 0 & 0 & 0 & s_{11} + s_{12} & s_{15} - s_{16} & s_{13} - s_{14}  \\
0 & b & 0 &  \beta_1 c \mu & 0 & 0 & 0 & 0 \\
0 & 0 & a &  0 & \beta_2 \mu & 0 & 0 & 0 \\
0 & - \sqrt{b} & 0 & - \mu & 0 & s_{51} + s_{52} & s_{55} - s_{56} & s_{53} - s_{54}  \\
0 &  \sqrt{b} & 0 & - \mu & 0 & s_{61} + s_{62} & s_{65} - s_{66} & s_{63} - s_{64}  \\
0 & 0 & - \sqrt{a} & 0 & - \mu & s_{31} + s_{32} & s_{35} - s_{36} & s_{33} - s_{34}  \\
0 & 0 &  \sqrt{a} & 0 & - \mu & s_{41} + s_{42} & s_{45} - s_{46} & s_{43} - s_{44}  \\
-1 & 0 & 0 & 0 & 0 & s_{21} + s_{22} & s_{25} - s_{26} & s_{23} - s_{24}  
\end{array} \right)
\end{equation}

\noindent where $s_{ij}$ means $S(\mu)_{ij}$ (cf. (\ref{Sdef}) for the definition of the matrix $S$. The dependency of $s_{ij}$ on $\mu$ is not always written in the following for the sake of simplicity). \\
\\
\noindent Finding eigenvalues and eigenvectors is possible if, and only if the system (\ref{syschareq}) has a non-vanishing solution, which is equivalent to the characteristic equation:

\begin{equation} \label{chareq}
\det[H(\mu)] = 0.
\end{equation} 

\noindent The computation of this determinant is realized through a triangularization. Using the classical Gauss pivot method leads to the equivalent characteristic equation:

\begin{equation} \label{chareq2}
a b (\beta_1 c + \sqrt{b})(\beta_2 + \sqrt{a}) \alpha_6 (\alpha_6 \gamma_7 - \gamma_6 \alpha_7) \tau = 0.
\end{equation}

\noindent where $\alpha_6, \alpha_7, \alpha_8, \gamma_6, \gamma_7, \gamma_8, \tau$ are defined by:

\begin{equation} \label{def_alpha_gamma} 
\left \{ \begin{array}{lllllll} 
\alpha_6:= -(\beta_1 c + \sqrt{b}) \sqrt{b} (s_{51} + s_{52}) - (\beta_1 c - \sqrt{b}) \sqrt{b} (s_{61} + s_{62}) \\
\alpha_7:= -(\beta_1 c + \sqrt{b}) \sqrt{b} (s_{55} - s_{56}) - (\beta_1 c - \sqrt{b}) \sqrt{b} (s_{65} - s_{66}) \\
\alpha_8:= -(\beta_1 c + \sqrt{b}) \sqrt{b} (s_{53} - s_{54}) - (\beta_1 c - \sqrt{b}) \sqrt{b} (s_{63} - s_{64}) \\
\gamma_6:= -(\beta_2 + \sqrt{a}) \sqrt{a} (s_{31} + s_{32}) - (\beta_2 - \sqrt{a}) \sqrt{b} (s_{41} + s_{42}) \\
\gamma_7:= -(\beta_2 + \sqrt{a}) \sqrt{a} (s_{35} - s_{36}) - (\beta_2 - \sqrt{a}) \sqrt{b} (s_{45} - s_{46}) \\
\gamma_8:= -(\beta_2 + \sqrt{a}) \sqrt{a} (s_{33} - s_{34}) - (\beta_2 - \sqrt{a}) \sqrt{b} (s_{43} - s_{44})  
\end{array} \right.
\end{equation}

\begin{equation} \label{def_tau} 
\tau:= \det \left( \begin{array}{cc} 
\alpha_6 \gamma_7 - \gamma_6 \alpha_7 & \alpha_6 \gamma_7 - \gamma_6 \alpha_7 \\
\alpha_6 (h_{27} - h_{17}) - \alpha_7 (h_{26} - h_{16}) & \alpha_6 (h_{28} - h_{18}) - \alpha_8 (h_{26} - h_{16})
\end{array} \right)
\end{equation}

\noindent with $h_{16}:= s_{11} + s_{12}$, $h_{26}:= s_{21} + s_{22}$, $h_{17}:= s_{15} - s_{16}$ and $h_{27}:= s_{25} - s_{26}$. \\
\\
The characteristic equation (\ref{chareq2}) gives three families of eigenvalues. The first one is the family of $\mu$'s satisfying $\alpha_6(\mu) =0$ or equivalently

\begin{equation} \label{1fam}
(\beta_1 c + \sqrt{b}) (s_{51}(\mu) + s_{52}(\mu)) + (\beta_1 c - \sqrt{b}) (s_{61}(\mu) + s_{62}(\mu)) = 0
\end{equation}

\noindent Now, to get the asymptotic behaviour of the large values of $\mu$, the $s_{ij}$'s are written as: 

\begin{equation} \label{sij_asymp} 
s_{ij}(\mu):= s^0_{ij} (e^{\mu D})_{ij} + \dfrac{(e^{\mu D})_{ij}}{\mu} s^{-1}_{ij} + o \left( \dfrac{1}{\mu} \right)
\end{equation}
%%%%%%%%%%%%%%%%%%%%%%%%%%%%%%%%%%%%%%%%%%%%%%%%%%%%%%%%%%%%%%%%%%%%%%%%%%%%%%%%%%%%%%%%%%%%%
\noindent The next step is to compute some of the terms $s^0_{ij}$ and $s^{-1}_{ij}$. \\
\\
\noindent From (\ref{Weq}), (\ref{decompoM}), the definition of $Z$ ($Z(\mu,x):=P^{-1} W(\mu, x)$) and (\ref{ZtoS2}), it follows:

\begin{equation} \label{deriveZ} 
Z_x(\mu,x) = \left[\mu D + P^{-1} M_0 P + \dfrac{1}{\mu} P^{-1} M_{-1} P \right] \sum_{k \geq 0} \left( \dfrac{S_{-k}(x)}{\mu^k} \right) \exp(\mu D x) Z(\mu,0). 
\end{equation}

\noindent Deriving (\ref{ZtoS2}) leads to:

\begin{equation} \label{deriveZ2} 
Z_x(\mu,x) = \sum_{k \geq 0} \left( \dfrac{S'_{-k}(x)}{\mu^k} + \dfrac{S_{-k}(x)}{\mu^k} \mu D \right) \exp(\mu D x) Z(\mu,0). 
\end{equation}

\noindent Identifying the terms of both sums (\ref{deriveZ}) and (\ref{deriveZ2}) gives:

\begin{equation} \label{S_calc} 
\left \{ \begin{array}{lll} 
S_0 D = D S_0 \\
S'_0 + S_{-1} D = D S_{-1} + P^{-1} M_0 P S_0 \\
\vdots \\
S'_{-k} + S_{-k-1} D = D S_{-k-1} + P^{-1} M_0 P S_{-k} + P^{-1} M_{-1} P S_{-k+1}, \, \forall \; k \geq 1. 
\end{array} \right.
\end{equation}
   
\noindent The dependency of $S_j$ on $x$ is not written here. \\      
\\
\noindent Assuming (\ref{cond}), the second item of Lemma \ref{lemD} applied to the first equation of (\ref{S_calc}), leads to: $S_0(x)$ is a diagonal matrix and so is its derivative $S'_0(x)$. A direct computation gives:

\begin{equation} \label{PM0}
P^{-1} M_0 P = \left( \begin{array}{cccccc} 
0 & 0 & -\frac{(a+1)l}{2\sqrt{a}} & \frac{(a+1)l}{2\sqrt{a}} & - \frac{1}{2 \sqrt{b}} & \frac{1}{2 \sqrt{b}} \\
0 & 0 & -\frac{(a+1)l}{2\sqrt{a}} & \frac{(a+1)l}{2\sqrt{a}} & - \frac{1}{2 \sqrt{b}} & \frac{1}{2 \sqrt{b}} \\
\frac{(a+1)l}{2\sqrt{a}} & \frac{(a+1)l}{2\sqrt{a}} & 0 & 0 & 0 & 0 \\
-\frac{(a+1)l}{2\sqrt{a}} & -\frac{(a+1)l}{2\sqrt{a}} & 0 & 0 & 0 & 0 \\
\frac{c}{2 \sqrt{b}} & \frac{c}{2 \sqrt{b}} & 0 & 0 & 0 & 0 \\
-\frac{c}{2 \sqrt{b}} & -\frac{c}{2 \sqrt{b}} & 0 & 0 & 0 & 0 
\end{array} \right)
\end{equation} 

\noindent From the first item of Lemma \ref{lemD} follows that the diagonal terms of $D S_{-1} - S_{-1} D$ vanish. So do the diagonal terms of $P^{-1} M_0 P$ (direct computation). \\
Thus, still assuming (\ref{cond}), the second equation of (\ref{S_calc}) implies: $S'_0(x) = 0, \forall x \in [0,1]$. \\
Then the function $S_0$ is constant and equal to the identity since $S_0(0)= I$. \\
This reduces the second equation of (\ref{S_calc}) to:

%%%%%%%%%%%%%%%%%%%%%%%%%%%%%%%%%%%%%%%%%%%%%%%%%%%%%%%%%%%%%%%%%%%
\begin{equation} \label{Sminus1}
S_{-1}(x) D - D S_{-1}(x) = P^{-1} M_0 P  
\end{equation}
%%%%%%%%%%%%%%%%%%%%%%%%%%%%%%%%%%%%%%%%%%%%%%%%%%%%%%%%%%%%%%%%%%%

\noindent This equation gives the non-diagonal terms of the matrix $S_{-1}(x)$ which do not depend on $x$, since the right-hand side of (\ref{Sminus1}) is constant with respect to $x$. \\
\\
Recall that, from (\ref{Sdef}) and (\ref{sij_asymp}) follows, for any $(i,j) \in [1,6]^2$ 

\begin{equation} \label{sij_calc} 
s^0_{ij} = S_0(1)_{ij} = \delta_{ij} \; \mbox{and} \; s^{-1}_{ij} = S_{-1}(1)_{ij}.
\end{equation}

%%%%%%%%%%%%%%%%%%%%  Lemma Calculus of S_minus_1 %%%%%%%%%%%%%%%%%%%%%%%%%%%%%%%%%%%%%%

\begin{lemma} (Non-diagonal terms of the matrix $S_{-1}$) \\ \label{lemS_minus_1}

\begin{itemize}
\item $s^{-1}_{12}=s^{-1}_{21}=s^{-1}_{34}=s^{-1}_{43}=s^{-1}_{35}=s^{-1}_{53}=s^{-1}_{36}=s^{-1}_{63}=s^{-1}_{45} = s^{-1}_{54}= s^{-1}_{46} = s^{-1}_{64}=s^{-1}_{56}=s^{-1}_{65} = 0.$

\item $s^{-1}_{13} = s^{-1}_{31} = s^{-1}_{42} = \dfrac{(a+1)l}{2(\sqrt{a}-1)}$, $s^{-1}_{14} = s^{-1}_{41} = s^{-1}_{23} = s^{-1}_{32} = -\dfrac{(a+1)l}{2(\sqrt{a}+1)}$, $s^{-1}_{24} = \dfrac{(a+1)l}{2(1-\sqrt{a})}$ \\
 $s^{-1}_{15} = \dfrac{1}{2(\sqrt{b}-1)}$, $s^{-1}_{26} = \dfrac{1}{2(1-\sqrt{b})}$, $s^{-1}_{16} = s^{-1}_{25} = -\dfrac{1}{2(\sqrt{b}+1)}$ \\
$ s^{-1}_{51} = s^{-1}_{62} = \dfrac{c}{2(\sqrt{b}-1)}$, $s^{-1}_{52} = s^{-1}_{61} = -\dfrac{c}{2(\sqrt{b}+1)}$.
\end{itemize}

\end{lemma}

\proof The non-diagonal terms of $S_{-1}$ are computed using (\ref{Sminus1}). \fin



%%%%%%%%%%%%%%%%%%%  Lemma Asymptotic behaviour of the first family of eigenvalues   %%%%%%%%%%

\begin{lemma} (Asymptotic behaviour of the first family of eigenvalues) \\ \label{asymp_1fam}

\noindent A first family of eigenvalues for the operator $\caA$ satisfies the characteristic equation:

\begin{equation} \label{1fam2}
(\beta_1 c + \sqrt{b}) (s_{51}(\mu) + s_{52}(\mu)) + (\beta_1 c - \sqrt{b}) (s_{61}(\mu) + s_{62}(\mu)) = 0. 
\end{equation}

\noindent This family satisfies $\mu_n = i n \pi + O(n^{-1})$, $n \in \Z$.

\end{lemma}

\proof Lemma \ref{lemS_minus_1} gives the asymptotic behaviour of $s_{51}(\mu)$, $s_{52}(\mu)$, $s_{61}(\mu)$ and $s_{62}(\mu)$. Inserting them in (\ref{1fam2}) leads to: $e^{\mu} - e^{- \mu} =0$. Hence the result.  \fin

%%%%%%%%%%%%%%%%%%%%%%%%%%%%%%%%%%%%%%%%%%%%%%%%%%%%%%%%%%%%%%%%%%%%%%%%%%%%%%%%%%%%%%%%%%%
\noindent This ends the proof of the Theorem. \fin \\
\\
%%%%%%%%%%%%%%%%%%%%%%%%%%%%%%%%%%%%%%%%%%%%%%%%%%%%%%%%%%%%%%%%%%%%%%%%%%%%%%%%%%%%%%%%%%%%%%%
\brk (Equal wave speeds) Note that, even if $(a,b)=(1,1)$ or $a=1$ with $b \neq 1$ or $b=1$ with $a \neq 1$ or $a=b$ with $a \neq 1$, the matrix $M_1$ in the proof of Theorem \ref{theig} is still diagonalizable with the same diagonalization as given at the beginning of the proof.  
\erk 
%%%%%%%%%%%%%%%%%%%%%%%%%%%%%%%%%%%%%%%%%%%%%%%%%%%%%%%%%%%%%%%%%%%%%%%%%%%%%%%%%%%%%%%%%%%%%%%
\brk (The inverse matrix $P^{-1}$) For the sake of completeness, the inverse matrix $P^{-1}$, involved in the diagonalization of the matrix $M_1$ in the proof of Theorem \ref{theig}, is the following one:
\begin{equation} \label{DefinvP}
P^{-1} = \left( \begin{array}{cccccc} 
1/2 & 0 & 0 & 1/2 & 0 & 0 \\
-1/2 & 0 & 0 & 1/2 & 0 & 0  \\
0 & 0 & 1/2 & 0 & 0 & \sqrt{a}/2 \\
0 & 0 & 1/2 & 0 & 0 & -\sqrt{a}/2 \\
0 & 1/2 & 0 & 0 & \sqrt{b}/2 & 0 \\
0 & 1/2 & 0 & 0 & -\sqrt{b}/2 & 0 
\end{array} \right).
\end{equation}
\erk 
%%%%%%%%%%%%%%%%%%%%%%%%%%%%%%%%%%%%%%%%%%%%%%%%%%%%%%%%%%%%%%%%%%%%%%%%%%%%%%%%%%%%%%%%%%%%
\brk \label{rkequalvelocities} (What happens if $(a,b)=(1,1)$ or $a=1, b \neq 1$ or $a \neq 1, b = 1$ or $a=b, a \neq 1$?) \\
If some of the wave velocities coincide, it is not possible anymore to conclude that $S_0(x)$ is diagonal. The terms $(S_0(x))_{13}$, $(S_0(x))_{31}$, $(S_0(x))_{24}$ and $(S_0(x))_{42}$ may not vanish and no information on them is available. \\
Only some diagonal terms of $P^{-1} M_0 P S_0$ are sure to vanish: the terms on the fifth and sixth lines. Thus the computation of the matrix $S_{-1}$ can not be done entirely: the terms $(S_{-1})_{11}$, $(S_{-1})_{22}$, $(S_{-1})_{33}$, $(S_{-1})_{44}$, $(S_{-1})_{13}$, $(S_{-1})_{31}$, $(S_{-1})_{24}$ and $(S_{-1})_{42}$ can not be computed. \\
As it is said in the introduction, the case when some wave velocities coincide can be treated by multiplier techniques and we are not interested in them in this paper.
\erk

%%%%%%%%%%%%%%%%%%%%%%%%%%%%%%%%%%%%%%%%%%%%%%%%
\brk \label{rkHanXu} (The paper by Z-J. Han, G-Q Xu, \cite{HanXu}) \\
The method we use here comes from (\cite{HanXu}). Their problem is more complicated: it is a thermoelastic Bresse system with variable coefficients. \\
The authors conclude that $S_0(x)$ is always diagonal. This is not true in general. They have omitted to exclude the case when the eigenvalues of their matrix $A_1$ are not simple. Even with constant coefficients, the assumption they should make is not clear because of the thermal effects. Our case is simpler since the matrix is diagonalized explicitly and the assumption we must make to be able to use Lemma \ref{lemD} is not physically relevant. 
\erk
%%%%%%%%%%%%%%%%%%%%%%%%%%%%%%%%%%%%%%%%%%%%%%%%%
\noindent The conclusion of this section is that uniform stability (i.e. exponential stability) does not hold since a subsequence of eigenvalues of the operator $\mathcal{A}$ is close to the imaginary axis. 

%%%%%%%%%%%%%%%%%%%%%%%%%%%%%%%%%%%%%%%%%%%%%%%%%%%%%%%%%%%%%%%%%%%%%%%%%%%%%%%%%%%%%%%%%%%%%%
%%%%%%%%%%%%%%%%%%%%%%%%%%%%%%%%%%%%%%%%%%%%%%%%%%%%%%%%%%%%%%%%%%%%%%%%%%%%%%%%%%%%%%%%%%%%%%%%%%%%%%%%%%%%%%%%%%%%%%%%%%%%%%%%%%%%%%%%%%%%%%
\section{Eigenfunctions of the operator $\caA$} \label{sec4}
%%%%%%%%%%%%%%%%%%%%%%%%%%%%%%%%%%%%%%%%%%%%%%%%%%%%%%%%%%%%%%%%%%%%%%%%%

\noindent In this section, we are interested in the asymptotic behaviour of the eigenfunctions, which can be deduced from the computations of the previous section. The following result could be useful to get the exact decay rate of the energy but a lot more work has to be done!

\bthe \label{theigfunc}(Asymptotic behaviour of the eigenfunctions)   \\
Assume Condition (\ref{condnorm}) as well as Condition (\ref{cond}) on $a$ and $b$. Define the functions $\varphi^0$, $\psi^0$ and $\omega^0$ by: 

\be \label{Phi0}
\left \{ \begin{array}{lll}
\varphi^0(\mu, x) := \dfrac{1}{\mu} \exp \left[- \mu \max \left \{ 2 ; 1 + \dfrac{1}{\sqrt{a}} ; 1 + \dfrac{1}{\sqrt{b}} \right \} \right] \exp \left[ \dfrac{\mu}{\tanh(\mu)} x \right] \\
\psi^0(\mu, x):= \dfrac{1}{\mu} \exp \left[- \mu \max \left \{ 2 ; 1 + \dfrac{1}{\sqrt{a}} ; 1 + \dfrac{1}{\sqrt{b}} \right \} \right] \exp \left[ \dfrac{\mu}{\sqrt{b} \tanh(\mu/\sqrt{b})} x \right] \\
\omega^0(\mu, x) := \dfrac{1}{\mu} \exp \left[- \mu \max \left \{ 2 ; 1 + \dfrac{1}{\sqrt{a}} ; 1 + \dfrac{1}{\sqrt{b}} \right \} \right] \exp \left[ \dfrac{\mu \tanh(\mu/\sqrt{a})}{\sqrt{a}} x \right] 
\end{array} \right.
\ee

\noindent for $\mu \in \C$ an eigenvalue of the operator $\mathcal{A}$, defined by (\ref{defDA}) and (\ref{defA}), and for $x \in [0,1]$. \\
\\
\noindent Then, if the function $\Phi(\mu,\cdot)$ defined by $\Phi(\mu,\cdot):= (\varphi(\mu,\cdot), \mu \varphi(\mu,\cdot), \psi(\mu,\cdot), \mu \psi(\mu,\cdot), \omega(\mu,\cdot), \mu \omega(\mu,\cdot))^t$ is an eigenfunction of the operator $\mathcal{A}$ associated with the eigenvalue $\mu$, it holds:

\be \label{asympPhi}
\left \{ \begin{array}{lll}
\dfrac{1}{\mu} \exp \left[- \mu \max \left \{ 2 ; 1 + \dfrac{1}{\sqrt{a}} ; 1 + \dfrac{1}{\sqrt{b}} \right \} \right] \varphi(\mu, x) = \varphi^0(\mu, x) + O \left( \dfrac{1}{\mu} \right) \\
\dfrac{1}{\mu} \exp \left[- \mu \max \left \{ 2 ; 1 + \dfrac{1}{\sqrt{a}} ; 1 + \dfrac{1}{\sqrt{b}} \right \} \right] \psi(\mu, x) = \psi^0(\mu, x) + O \left( \dfrac{1}{\mu} \right) \\
\dfrac{1}{\mu} \exp \left[- \mu \max \left \{ 2 ; 1 + \dfrac{1}{\sqrt{a}} ; 1 + \dfrac{1}{\sqrt{b}} \right \} \right] \omega(\mu, x) = \omega^0(\mu, x) + O \left( \dfrac{1}{\mu} \right) 
\end{array} \right.
\ee

\noindent and $\Phi^0(\mu,\cdot)$ defined by $\Phi^0(\mu,\cdot):= (\varphi^0(\mu,\cdot), \mu \varphi^0(\mu,\cdot), \psi^0(\mu,\cdot), \mu \psi^0(\mu,\cdot), \omega^0(\mu,\cdot), \mu \omega^0(\mu,\cdot))^t$ satisfies:  
 
$$\| \Phi^0(\mu,\cdot) \|_{\mathcal{H}} = O(1).$$ 

\ethe

\proof Denote by $\Phi(\mu,x):= (\varphi(\mu,x), \mu \varphi(\mu,x), \psi(\mu,x), \mu \psi(\mu,x), \omega(\mu,x), \mu \omega(\mu,x))^t$ an eigenfunction of the operator $\mathcal{A}$ defined by (\ref{defDA}) and (\ref{defA}). \\
Starting from the third equation of (\ref{ZtoS2}) where $Z$ is defined by (\ref{defZ}), it holds:

\be
W(\mu, x) = P  \Sigma(\mu,x) \exp(\mu D x) P^{-1} W(\mu,0)
\ee

\noindent where the vector $W(\mu,x)$ is given by (\ref{defW}), $\Sigma(\mu,x) := \sum_{k \geq 0} \dfrac{S_{-k}(x)}{\mu^k}$ and $S_{-k}(x)$ is a square matrix of order $6$ which can be computed step by step as in the proof of Theorem \ref{theig}. The matrices $P$, $D$ and $P^{-1}$ are defined by (\ref{PdefD}) and (\ref{DefinvP}). \\
A direct calculation leads to: $W(\mu, x) = $

\be
P  \Sigma(\mu,x) \dfrac{1}{2 \mu} \left(e^{\mu} \varphi_x(\mu,0), e^{-\mu} \varphi_x(\mu,0), e^{\frac{\mu}{\sqrt{a}}} \sqrt{a} \omega_x(\mu,0), e^{-\frac{\mu}{\sqrt{a}}} \sqrt{a} \omega_x(\mu,0), e^{\frac{\mu}{\sqrt{b}}} \sqrt{b} \psi_x(\mu,0),- e^{-\frac{\mu}{\sqrt{b}}} \sqrt{b} \psi_x(\mu,0) \right)^t
\ee

\begin{lemma} (Asymptotic behaviour of $\Sigma(\mu,x)$ with respect to $\mu$) \\
The asymptotic behaviour of $\Sigma(\mu,x)$ is given by

$$[\Sigma(\mu,x)]_{ij} = \left \{ \begin{array}{ll}
1 + \dfrac{c_k x}{\mu} + O \left( \dfrac{1}{\mu^2} \right) \; \mbox{if} \; i = j = k\\
\dfrac{s_{ij}^{-1}}{\mu} + O \left( \dfrac{1}{\mu^2} \right) \; \mbox{if} \; i \neq j
\end{array} \right. $$

\noindent where the values of $s_{ij}^{-1}$ are computed in Lemma \ref{lemS_minus_1}. 

\end{lemma}

\proof By definition, $[\Sigma(\mu,x)]_{ij} = [S_0(x)]_{ij} + \dfrac{1}{\mu} [S_{-1}(x)]_{ij} + O \left( \dfrac{1}{\mu^2} \right)$. In the proof of Theorem \ref{theig}, the matrix $S_0(x)$ was proved to be equal to the identity matrix and the non-diagonal terms of the matrix $S_{-1}(x)$ are proved to be independent of $x$. There only remains to prove that the diagonal terms of $S_{-1}(x)$ vanish. \\
Recall the third equation of (\ref{S_calc}) where $M_0$, $M_1$ and $M_{-1}$ are defined by (\ref{defM1}) and (\ref{defMmoins1}):

\begin{equation} 
S'_{-1} + \underbrace{S_{-2} D - D S_{-2}}_{\mbox{the diagonal terms vanish}} = \underbrace{P^{-1} M_0 P S_{-1}}_{\mbox{independent of} \, x} + \underbrace{P^{-1} M_{-1} P S_0}_{\mbox{independent of} \, x}.
\end{equation}   

\noindent The properties follow from the fact that the diagonal terms of $P^{-1} M_0 P$ vanish (cf. (\ref{PM0})) and the non-diagonal terms of $S_{-1}(x)$ are independent of $x$. \\
Thus, there exist constants $c_k$ such that: 

\begin{equation} 
[S'_{-1}(x)]_{kk} = c_k.
\end{equation} 

\noindent And, since $S_{-1}(0) = 0$ (it is a consequence of (\ref{ZtoS2})), it holds $[S_{-1}(x)]_{kk} = c_k x$. \fin \\
\\
\noindent Define by $\tilde{W}(\mu, x)$ the first term of the asymptotic expansion of $W(\mu, x)$:

\be 
\begin{array}{ll}
\tilde{W}(\mu, x) = P I \dfrac{1}{2 \mu} \left(e^{\mu} \tilde{\varphi}_x(\mu,0), e^{-\mu} \tilde{\varphi}_x(\mu,0), e^{\frac{\mu}{\sqrt{a}}} \sqrt{a} \omega_x(\mu,0), e^{-\frac{\mu}{\sqrt{a}}} \sqrt{a} \omega_x(\mu,0), e^{\frac{\mu}{\sqrt{b}}} \sqrt{b} \tilde{\psi}_x(\mu,0),- e^{-\frac{\mu}{\sqrt{b}}} \sqrt{b} \tilde{\psi}_x(\mu,0) \right)^t \\
= \dfrac{1}{2 \mu} \left(
\begin{array}{llllll}
e^{\mu} \tilde{\varphi}_x(\mu,0) - e^{-\mu} \tilde{\varphi}_x(\mu,0) \\
e^{\frac{\mu}{\sqrt{b}}} \sqrt{b} \tilde{\psi}_x(\mu,0) - e^{-\frac{\mu}{\sqrt{b}}} \sqrt{b} \tilde{\psi}_x(\mu,0) \\
e^{\frac{\mu}{\sqrt{a}}} \sqrt{a} \tilde{\omega}_x(\mu,0) + e^{-\frac{\mu}{\sqrt{a}}} \sqrt{a} \tilde{\omega}_x(\mu,0) \\
e^{\mu} \tilde{\varphi}_x(\mu,0) + e^{-\mu} \tilde{\varphi}_x(\mu,0) \\
e^{\frac{\mu}{\sqrt{b}}} \sqrt{b} \tilde{\psi}_x(\mu,0) + e^{-\frac{\mu}{\sqrt{b}}} \sqrt{b} \tilde{\psi}_x(\mu,0) \\
e^{\frac{\mu}{\sqrt{a}}} \sqrt{a} \tilde{\omega}_x(\mu,0) - e^{-\frac{\mu}{\sqrt{a}}} \sqrt{a} \tilde{\omega}_x(\mu,0) \end{array} \right).
\end{array}
\ee

\noindent By definition of $\tilde{W}$ (see (\ref{defW}) mutatis mutandis), the following system follows:

\be 
\left \{ \begin{array}{llllll}
\tilde{\varphi}(\mu, x) = \dfrac{1}{2 \mu} (e^{\mu} - e^{-\mu}) \tilde{\varphi}_x(\mu, 0) \\
\tilde{\psi}(\mu, x) = \dfrac{1}{2 \mu}  \sqrt{b} (e^{\frac{\mu}{\sqrt{b}}} - e^{-\frac{\mu}{\sqrt{b}}}) \tilde{\psi}_x(\mu, 0) \\
\tilde{\omega}(\mu, x) = \dfrac{1}{2 \mu} \sqrt{a} (e^{\frac{\mu}{\sqrt{a}}} + e^{-\frac{\mu}{\sqrt{a}}}) \tilde{\omega}_x(\mu, 0) \\
\dfrac{\tilde{\varphi}_x(\mu, x)}{\mu} = \dfrac{1}{2 \mu} (e^{\mu} + e^{-\mu}) \tilde{\varphi}_x(\mu, 0) \\
\dfrac{\tilde{\psi}_x(\mu, x)}{\mu} = \dfrac{1}{2 \mu}  \sqrt{b} (e^{\frac{\mu}{\sqrt{b}}} + e^{-\frac{\mu}{\sqrt{b}}}) \tilde{\psi}_x(\mu, 0) \\
\dfrac{\tilde{\omega}_x(\mu, x)}{\mu} = \dfrac{1}{2 \mu} \sqrt{a} (e^{\frac{\mu}{\sqrt{a}}} - e^{-\frac{\mu}{\sqrt{a}}}) \tilde{\omega}_x(\mu, 0) 
\end{array} \right.
\ee

\noindent Since $\mu \neq i n \pi, n \in \Z$, $\mu \neq i n \sqrt{a} \dfrac{\pi}{2}, n \in \Z$ and $\mu \neq i n \sqrt{b} \pi, n \in \Z$ (there is no eigenvalue on the imaginary axis, cf. the proof of Theorem \ref{thstrong}), it holds: 

\be 
\left \{ \begin{array}{lll}
\dfrac{\tilde{\varphi}_x(\mu, x)}{\tilde{\varphi}(\mu, x)} = \mu \dfrac{e^{\mu} + e^{-\mu}}{e^{\mu} - e^{-\mu}} \\
\dfrac{\tilde{\psi}_x(\mu,x)}{\tilde{\psi}(\mu, x)} = \mu \dfrac{e^{\frac{\mu}{\sqrt{b}}} + e^{-\frac{\mu}{\sqrt{b}}}}{\sqrt{b} \left(e^{\frac{\mu}{\sqrt{b}}} - e^{-\frac{\mu}{\sqrt{b}}} \right)} \\
\dfrac{\tilde{\omega}_x(\mu,x)}{\tilde{\omega}(\mu, x)} = \mu \dfrac{e^{\frac{\mu}{\sqrt{a}}} - e^{-\frac{\mu}{\sqrt{a}}}}{\sqrt{a} \left(e^{\frac{\mu}{\sqrt{a}}} + e^{-\frac{\mu}{\sqrt{a}}} \right)}.
\end{array} \right.
\ee

\noindent One solution of this system is:

\be 
\left \{ \begin{array}{lll}
\tilde{\varphi}(\mu, x) = \exp \left[ \dfrac{\mu}{\tanh(\mu)} x \right] \\
\tilde{\psi}(\mu, x) = \exp \left[ \dfrac{\mu}{\sqrt{b} \tanh(\mu/\sqrt{b})} x \right] \\
\tilde{\omega}(\mu, x) = \exp \left[ \dfrac{\mu \tanh(\mu/\sqrt{a})}{\sqrt{a}} x \right] 
\end{array} \right.
\ee

\noindent The computation of $\| \tilde{\Phi}(\mu,\cdot) \|_{\mathcal{H}}$ where $ \tilde{\Phi}(\mu,\cdot):= (\tilde{\varphi}(\mu,\cdot), \mu \tilde{\varphi}(\mu,\cdot), \tilde{\psi}(\mu,\cdot), \mu \tilde{\psi}(\mu,\cdot), \tilde{\omega}(\mu,\cdot), \mu \tilde{\omega}(\mu,\cdot))^t \|$ leads to the choice of the functions: 

\be \label{Phi0}
\left \{ \begin{array}{lll}
\varphi^0(\mu, x) = \dfrac{1}{\mu} \exp \left[- \mu \max \left \{ 2 ; 1 + \dfrac{1}{\sqrt{a}} ; 1 + \dfrac{1}{\sqrt{b}} \right \} \right] \exp \left[ \dfrac{\mu}{\tanh(\mu)} x \right] \\
\psi^0(\mu, x)= \dfrac{1}{\mu} \exp \left[- \mu \max \left \{ 2 ; 1 + \dfrac{1}{\sqrt{a}} ; 1 + \dfrac{1}{\sqrt{b}} \right \} \right] \exp \left[ \dfrac{\mu}{\sqrt{b} \tanh(\mu/\sqrt{b})} x \right] \\
\omega^0(\mu, x) = \dfrac{1}{\mu} \exp \left[- \mu \max \left \{ 2 ; 1 + \dfrac{1}{\sqrt{a}} ; 1 + \dfrac{1}{\sqrt{b}} \right \} \right] \exp \left[ \dfrac{\mu \tanh(\mu/\sqrt{a})}{\sqrt{a}} x \right] 
\end{array} \right.
\ee

\noindent Thus $\| \Phi^0(\mu,\cdot) \|_{\mathcal{H}}= O(1)$ where $\Phi^0(\mu,\cdot) := (\varphi^0(\mu,\cdot), \mu \varphi^0(\mu,\cdot), \psi^0(\mu,\cdot), \mu \psi^0(\mu,\cdot), \omega^0(\mu,\cdot), \mu \omega^0(\mu,\cdot))^t \|$.  \fin \\




%For acknowledgements section, please don't number the section, please begin it with \section*{Acknowledgements}
%\section*{Acknowledgments} The authors thank the referees for their attentive reading of the manuscript and their valuable remarks.

% You may incorporate your references as follows in your main tex file.
% Using BibTex is not recommended but can be handled.

\begin{thebibliography}{99}

\bibitem{AkilBad} \newblock M. Akil and H. Badawi, \newblock {The influence of the physical coefficients of a Bresse system with one singular local viscous damping in the longitudinal displacement on its stabilization}, \emph{Evol. Equ. Control Theory}, \textbf{11(6)} (2022), 1903--1928.

\bibitem{AkilBadNicWeh} \newblock M. Akil, H. Badawi, S. Nicaise and A. Wehbe, \newblock{On the stability of Bresse system with one discontinuous local internal Kelvin-Voigt damping on the axial force} \emph{Z. Angew. Math. Phys.}, \textbf{72, No 3, Art. 126} (2021).

\bibitem{AkilChiGhaWeh} \newblock M. Akil, Y. Chitour, M. Ghader, and A. Wehbe. \newblock{Stability and exact controllability of a Timoshenko system with only one fractional damping on the boundary}, \emph{Asymptotic Analysis}, \textbf{119(3-4)} (2020), 221--280.

\bibitem{Ammar2} \newblock F. Ammar-Kodja, S. Kerbal and A. Soufyane,
\newblock {Stabilization of the nonuniform Timoshenko beam}, \emph{J. Math. Anal. Appl.}, \textbf{327} (2007), 525--538.

\bibitem{AVRR} \newblock M.S. Alves, J.E.  Mu\~{n}oz Rivera, O.P. Vera Villagran, A. Rambaud, 
\newblock{Exponential stability to the Bresse system with boundary dissipation conditions,} 
https://doi.org/10.48550/arXiv.1506.01657 (2015), 10 pp.

\bibitem{AB} W. Arendt, C.J.K. Batty, Tauberian theorems and stability of one-parameter of semi-groups, Trans. Amer. Math. Soc. \bf{305(2)} \rm (1988), 837-852.

\bibitem{mr1} \newblock M. Bassam, D. Mercier, S. Nicaise and A. Wehbe, 
\newblock{Polynomial stability of the Timoshenko system by one boundary damping,} \emph{J. Math. Anal and Appl.}, \textbf{425(2)} (2015), 1--27.

\bibitem{HanXu} \newblock Z-J. Han, G-Q Xu, \newblock{Spectral analysis and stability of thermoelastic Bresse system with second sound and boundary viscoelastic damping}, \emph{Math. Meth. Appl. Sc.} \textbf{38(1)} (2013), 94--112.
% https://doi.org/10.1002/mma.3052

\bibitem{KR} 
\newblock J. U. Kim and Y. Renardy,
\newblock {Boundary control of the Timoshenko beam}, \emph{SIAM J. Control Optim.}, \textbf{25(6)} (1987), 1417--1429. 

\bibitem{LiuRao} \newblock Z. Liu and B. Rao,
\newblock {Energy decay rate of the thermoelastic Bresse system}, \emph{Z. Angew. Math. Phys.}, \textbf{60} (2009), 54--69. 

\bibitem{EECT18} \newblock D. Mercier, V. R�gnier, \newblock{Decay rate of the Timoshenko system with one boundary damping}, \emph{Evol. Equ. Control Theory}, \textbf{8(2)} (2019), 423--445.

\bibitem{P} A. Pazy, Semigroups of Linear Operators and Applications to Partial Differential Equations, Applied Mathematical Sciences \bf{44} \rm(1983), Springer-Verlag.

\bibitem{RN} J. E. Mu�oz Rivera, M.G. Naso, \newblock{Boundary stabilization of Bresse systems}. \emph{Z. Angew. Math. Phys.} \bf{70,  No 2, Art. 56} \rm(2019) (2019), 16 pp. 

\bibitem{Z-R-M} \newblock L. Zietsman, N.F.J. van Rensburg and A.J. van der Merwe, 
\newblock{A Timoshenko beam with tip body and boundary damping}, \emph{Wave Motion}, \textbf{39} (2004), 199--211.

\bibitem{Santos} \newblock M.L. Santos, D.S.A. Junior, 
\newblock{Numerical Exponential Decay to Dissipative Bresse System}, \emph{J. Appl. Math.}, vol. 2010, Article ID 848620 (2010), 17 pp.



\end{thebibliography}
\end{document}











































